\section{Multiplicative functions: work one prime at a time}
\label{sec:multiplicative-functions}

The identity $\varphi(mn)=\varphi(m)\varphi(n)$ for coprime $m,n$
suggests a useful class of functions. Many counts become products when
their prime factors can be chosen independently.

\begin{definition}
An \emph{arithmetic function} assigns a real number $f(n)$ to each
positive integer $n$. It is \emph{multiplicative} if $f(1)=1$ and
\[
 f(mn)=f(m)f(n)\qquad\text{whenever }\gcd(m,n)=1.
\]
It is \emph{completely multiplicative} if this equality holds for all
positive $m,n$, without the coprimality condition.
\end{definition}

Unique factorization now gives
\[
 f\left(\prod_i p_i^{k_i}\right)=\prod_i f(p_i^{k_i}).
\]
Thus a multiplicative function is determined by its values at prime
\emph{powers}. A completely multiplicative function needs only its
values at primes, since then $f(p^k)=f(p)^k$.
The constant function $\mathbf{1}(n)=1$ and the identity function
$\operatorname{id}(n)=n$ are completely multiplicative.
Euler's totient is multiplicative by CRT, but
$\varphi(4)=2\ne1=\varphi(2)^2$.

\subsection*{Two more examples: count divisors and add them}

Write $\sum_{d\mid n}$ for a sum over the positive divisors of $n$, and set
\[
 \tau(n)=\sum_{d\mid n}1,\qquad \sigma(n)=\sum_{d\mid n}d.
\]
If $n=\prod_i p_i^{k_i}$, each divisor chooses an exponent between
$0$ and $k_i$ independently at each prime. Counting choices, or summing
the resulting products, gives
\[
 \tau(n)=\prod_i(k_i+1),\qquad
 \sigma(n)=\prod_i(1+p_i+\cdots+p_i^{k_i}).
\]
In particular, $\tau(p^k)=k+1$ and
$\sigma(p^k)=(p^{k+1}-1)/(p-1)$. Both functions are multiplicative:
coprime factors have disjoint prime choices. For $72=2^3\cdot3^2$,
\[
 \tau(72)=4\cdot3=12,\qquad
 \sigma(72)=(1+2+4+8)(1+3+9)=195.
\]

\subsection*{Divisor sums and Dirichlet convolution}

\begin{definition}
The \emph{Dirichlet convolution} of arithmetic functions $f,g$ is
\[
 (f*g)(n)=\sum_{d\mid n}f(d)g(n/d).
\]
Its identity is the function $\varepsilon$ given by
$\varepsilon(1)=1$ and $\varepsilon(n)=0$ for $n>1$.
\end{definition}

For example, $f*\mathbf{1}$ is the divisor-sum function
$n\mapsto\sum_{d\mid n}f(d)$, so
$\tau=\mathbf{1}*\mathbf{1}$ and $\sigma=\operatorname{id}*\mathbf{1}$.
Swapping $d$ and $n/d$ proves commutativity. Both ways to form a triple
convolution sum $f(a)g(b)h(c)$ over $abc=n$, proving associativity.
Only the term $d=n$ survives in $f*\varepsilon$, giving $f(n)$.
These are finite sums; no convergence questions enter.

\begin{proposition}
The convolution of two multiplicative functions is multiplicative.
In particular, divisor sums preserve multiplicativity.
\end{proposition}
\begin{proof}
For coprime $m,n$, each divisor of $mn$ has a unique form $ab$ with
$a\mid m$ and $b\mid n$. The pairs $a,b$ and $m/a,n/b$ are coprime.
Therefore
\[
 \begin{aligned}
 (f*g)(mn)
 &=\sum_{a\mid m}\sum_{b\mid n} f(a)f(b)g(m/a)g(n/b)\\
 &=\left(\sum_{a\mid m}f(a)g(m/a)\right)
   \left(\sum_{b\mid n}f(b)g(n/b)\right).
 \end{aligned}
\]
Also $(f*g)(1)=f(1)g(1)=1$.
\end{proof}

\begin{warningbox}[breakable=false]
Convolution is not pointwise multiplication. For example,
$(\mathbf{1}*\mathbf{1})(6)=4$, although $\mathbf{1}(6)\mathbf{1}(6)=1$.
And multiplicativity does not let us replace $f(p^k)$ by $f(p)^k$:
the coprimality hypothesis is doing real work.
\end{warningbox}

\subsection*{The totient identity to be inverted}

\begin{proposition}
For every positive integer $n$,
\[
 \sum_{d\mid n}\varphi(d)=n,
 \qquad\text{or equivalently}\qquad
 \varphi*\mathbf{1}=\operatorname{id}.
\]
\end{proposition}
\begin{proof}
Partition the integers $1\le r\le n$ according to $g=\gcd(r,n)$.
For each divisor $g\mid n$, put $d=n/g$ and write $r=gs$.
Then $1\le s\le d$ and $\gcd(s,d)=1$, so this class has
$\varphi(d)$ elements. Every $r$ belongs to exactly one class.
Adding their sizes counts all $n$ integers, including $r=n$
in the class corresponding to $d=1$.
\end{proof}

For $n=12$, this reads
$\varphi(1)+\varphi(2)+\varphi(3)+\varphi(4)+\varphi(6)+\varphi(12)
=1+1+2+2+2+4=12$.
Can we recover the individual values from their divisor sums?

\section{M\"obius inversion: undo a divisor sum}
\label{sec:mobius-inversion}

If $F(n)=\sum_{d\mid n}f(d)$, we can solve successively:
$f(1)=F(1)$ and $f(n)=F(n)-\sum_{d\mid n,\,d<n}f(d)$.
M\"obius inversion packages those repeated subtractions into one formula.
The M\"obius function introduced here is a different object from the
linear fractional transformations in Chapter~3.

\begin{definition}
An integer is \emph{squarefree} if no prime square divides it.
The \emph{M\"obius function} is
\[
 \mu(n)=
 \begin{cases}
 1,&n=1,\\
 (-1)^r,&n\text{ is a product of }r\text{ distinct primes},\\
 0,&\text{some }p^2\text{ divides }n.
 \end{cases}
\]
\end{definition}

For example, $\mu(6)=1$, $\mu(12)=0$, and $\mu(30)=-1$.
For coprime inputs, square factors remain square factors; if neither
input has one, their disjoint lists of prime factors concatenate.
This proves that $\mu$ is multiplicative. It is not completely
multiplicative: $\mu(p^2)=0$ but $\mu(p)^2=1$.

\begin{lemma}[Cancellation over the divisors]
For every positive integer $n$,
\[
 \sum_{d\mid n}\mu(d)=\varepsilon(n).
\]
\end{lemma}
\begin{proof}
For $n=1$, the sum is $\mu(1)=1$. If $n>1$, choose a prime $p\mid n$.
Only squarefree divisors contribute. Pair each squarefree divisor
not containing $p$ with the divisor obtained by adjoining $p$.
The two terms have opposite signs, and every nonzero term belongs
to exactly one pair. The sum is zero.
\end{proof}

\begin{theorem}[M\"obius inversion on divisors]
For arbitrary arithmetic functions $f,F$,
\[
 F(n)=\sum_{d\mid n}f(d)\quad\text{for every }n
 \quad\Longleftrightarrow\quad
 f(n)=\sum_{d\mid n}\mu(d)F(n/d)\quad\text{for every }n.
\]
Equivalently, $f(n)=\sum_{d\mid n}\mu(n/d)F(d)$.
Neither function needs to be multiplicative.
\end{theorem}
\begin{proof}
Substitute the first formula into the proposed inverse and collect
the coefficient of each $f(e)$:
\[
 \begin{aligned}
 \sum_{d\mid n}\mu(d)F(n/d)
 &=\sum_{d\mid n}\mu(d)\sum_{e\mid n/d}f(e)\\
 &=\sum_{e\mid n}f(e)\sum_{d\mid n/e}\mu(d)=f(n).
 \end{aligned}
\]
The inner sum vanishes unless $e=n$. In convolution notation the
cancellation lemma says $\mu*\mathbf{1}=\varepsilon$, so the reverse
implication follows as well:
$f=\mu*F$ gives $\mathbf{1}*f=(\mathbf{1}*\mu)*F=F$.
\end{proof}

\begin{example}[Recover the totient]
Invert $\sum_{d\mid n}\varphi(d)=n$ to obtain
\[
 \varphi(n)=\sum_{d\mid n}\mu(d)\frac nd
 =n\sum_{d\mid n}\frac{\mu(d)}d
 =n\prod_{p\mid n}\left(1-\frac1p\right).
\]
For the last equality, expand the finite product: choosing a subset
of the distinct prime factors produces exactly one squarefree divisor,
with its M\"obius sign. For $n=12$, only $d=1,2,3,6$ contribute:
\[
 \varphi(12)=12-6-4+2=4.
\]
CRT proved this product formula earlier; inversion explains the
alternating subtractions behind it.
\end{example}

\subsection*{Return to necklaces of any length}
\label{sec:primitive-necklaces}

Fix $a\ge1$ colors. Let $P_a(n)$ count labeled words of length $n$
whose \emph{least} rotational period is $n$; call them primitive words.
Every word of length $n$ has a unique least period $d\mid n$ and
is obtained by repeating a unique primitive block of length $d$.
Consequently
\[
 a^n=\sum_{d\mid n}P_a(d),\qquad
 P_a(n)=\sum_{d\mid n}\mu(d)a^{n/d}.
\]
The primitive words have rotation orbits of size $n$. Thus
\[
 \frac1n\sum_{d\mid n}\mu(d)a^{n/d}
\]
is the integer number of primitive necklaces. At $n=p^k$, only
$d=1,p$ survive, recovering $a^{p^k}-a^{p^{k-1}}$ and the Euler proof.
For binary words of length $6$,
\[
 P_2(6)=2^6-2^3-2^2+2=54,
\]
so there are nine primitive necklaces.

\begin{warningbox}[breakable=false]
This counts necklaces of \emph{full} period. To count all length-$n$
necklaces, include the shorter periods too:
\[
 N_a(n)=\sum_{d\mid n}\frac{P_a(d)}d.
\]
Indeed, a primitive block of length $d$ has $d$ distinct rotations
even after it is repeated to length $n$. For $a=2,n=6$, the contributions
are $2,1,2,9$ for $d=1,2,3,6$, giving $N_2(6)=14$.
\end{warningbox}

\begin{checkpoint}
Divisor sums record a total after grouping objects by a divisor.
M\"obius inversion recovers the exact groups. In the totient example
the grouping uses a gcd; in the necklace example it uses a least period.
The next section explains how the same method works for other ways of
ordering the groups. It is bonus material and can be read independently
of the later group-theory chapter.
\end{checkpoint}
